\section{Trigonometric integrals}

Trigonometric integrals require two kinds of reading. First, we should recognize direct derivatives, such as \((\sin x)'=\cos x\). Second, we should know a few identities that rewrite products or powers into simpler forms.

The goal in this section is not to list every possible trigonometric integral. The goal is to learn how to choose a useful transformation.

\subsection*{Direct substitution}

Some trigonometric integrals are simple substitutions.

\begin{examplebox}
Compute
\[
\int \sin^3 x\cos x\,dx.
\]
The inside expression is \(\sin x\), and its derivative is \(\cos x\). Let
\[
u=\sin x.
\]
Then
\[
du=\cos x\,dx.
\]
The integral becomes
\[
\int u^3\,du=\frac{u^4}{4}+C.
\]
Therefore
\[
\int \sin^3 x\cos x\,dx=\frac{\sin^4 x}{4}+C.
\]
\end{examplebox}

The integrand was not difficult after the inside function was recognized.

\subsection*{Using identities}

Products of sines and cosines may require identities. Two useful identities are
\[
\sin^2 x=\frac{1-\cos 2x}{2},
\qquad
\cos^2 x=\frac{1+\cos 2x}{2}.
\]
These are useful for even powers of sine or cosine.

\begin{examplebox}
Compute
\[
\int \sin^2 x\,dx.
\]
Use the identity
\[
\sin^2 x=\frac{1-\cos 2x}{2}.
\]
Then
\[
\int \sin^2 x\,dx
=\frac12\int 1\,dx-\frac12\int \cos 2x\,dx.
\]
Since
\[
\int \cos 2x\,dx=\frac12\sin 2x,
\]
we get
\[
\int \sin^2 x\,dx=\frac{x}{2}-\frac{\sin 2x}{4}+C.
\]
\end{examplebox}

The identity changed a power into a sum of simpler terms.

\subsection*{Products with different angles}

Products such as \(\sin ax\cos bx\) or \(\cos ax\cos bx\) can be simplified using product-to-sum identities:
\[
\sin A\cos B=\frac12\bigl(\sin(A+B)+\sin(A-B)\bigr),
\]
\[
\cos A\cos B=\frac12\bigl(\cos(A+B)+\cos(A-B)\bigr),
\]
\[
\sin A\sin B=\frac12\bigl(\cos(A-B)-\cos(A+B)\bigr).
\]

\begin{examplebox}
Compute
\[
\int \sin 3x\cos 2x\,dx.
\]
Use
\[
\sin A\cos B=\frac12(\sin(A+B)+\sin(A-B)).
\]
With \(A=3x\), \(B=2x\),
\[
\sin 3x\cos 2x=\frac12(\sin 5x+\sin x).
\]
Therefore
\[
\int \sin 3x\cos 2x\,dx
=\frac12\int \sin 5x\,dx+\frac12\int \sin x\,dx.
\]
Thus
\[
\int \sin 3x\cos 2x\,dx
=-\frac1{10}\cos 5x-\frac12\cos x+C.
\]
\end{examplebox}

\subsection*{Tangent and secant patterns}

The derivative of \(\tan x\) is
\[
\frac{1}{\cos^2 x}.
\]
Therefore expressions containing \(1/\cos^2 x\) often lead to tangent.

\begin{examplebox}
Compute
\[
\int \frac{dx}{2\cos^2 3x}.
\]
Rewrite:
\[
\int \frac{dx}{2\cos^2 3x}=\frac12\int \frac{1}{\cos^2 3x}\,dx.
\]
Let
\[
u=3x,
\qquad
du=3\,dx.
\]
Then
\[
\frac12\int \frac{1}{\cos^2 3x}\,dx
=\frac16\int \frac{1}{\cos^2 u}\,du.
\]
Thus
\[
\int \frac{dx}{2\cos^2 3x}=\frac16\tan(3x)+C.
\]
\end{examplebox}

\begin{summarybox}
When integrating trigonometric expressions, ask:
\begin{itemize}
\item Is there an inside function whose derivative is present?
\item Can substitution use \(\sin x\), \(\cos x\), or a multiple of \(x\)?
\item Are even powers present, suggesting half-angle identities?
\item Is there a product of different angles, suggesting product-to-sum identities?
\item Does the integrand contain \(1/\cos^2 x\), suggesting tangent?
\end{itemize}
\end{summarybox}

Trigonometric integration is largely the art of rewriting the expression until a familiar derivative appears backward.